\section{问题重述与分析}

\subsection{问题背景}

2026 年全国大学生数学建模竞赛 B 题要求在一半径为 $1800$ m 的圆形作业区域内，派遣一台
机器狗对可能散布在区域内的干扰源完成"搜索—定位—清除"全流程。区域内共 $20$ 个频道，
每个频道至多一个干扰源，源随局随机生成；机器狗可沿任意方向以 $5$ m/s 移动，在任意位置
对任意频道测向需 $5$ s（连续测两频道各需 $1$ s 的频道切换代价），对估计位置 $\le 5$ m
的源可直接近距清除。任务评判标准为\textbf{清除率 $100\%$ 且虚拟总耗时最短}。

\subsection{问题四与问题三的唯一差异}

前三问均为\textbf{全向干扰源}：任意位置、距源 $\le$ 接收半径即可听到，因此"听不到"
必然意味着"太远"。问题四引入\textbf{定向干扰源}：定向源只在\textbf{定向方向两侧各
$90^\circ$}（全角 $180^\circ$，含边界）的波束内有信号，波束之外一律返回
\texttt{no\_signal}（官方引擎判据：$|p-g|\le R$ 且在波束角内）。

\begin{quote}
\textbf{"搜索不到"不再等价于"源太远"，还可能是方向不对。}实测：对定向源正北 $300$ m
测得方向 $270.8^\circ$，正南 $300$ m 则 \texttt{no\_signal}，而正南 $19$ m 直接清除仍
\texttt{success}。
\end{quote}

据此，问题三的两处核心设计失效，须结构性重做：
\begin{itemize}
  \item 问题三的 $7$ 点覆盖只保证"圆域内任意点到最近巡视点 $\le 1000$ m"（对全向源必
        听到）；对定向源，\textbf{背对}所有巡视点的源即使近在咫尺也听不到（$200$ 万随机
        源蒙特卡洛实测：$7$ 点听到率仅 $80.7\%$，平均每 $5$ 个源漏 $1$ 个）。
  \item 问题三把 \texttt{no\_signal} 作为"源在检测点 $1000$ m 外"的定位硬约束；问题四中
        \texttt{no\_signal} 可能是方向不对，不能转成"圆盘外"约束，定位区域只能使用
        \texttt{direction} 与 \texttt{near} 两类约束（见第~\ref{sec:localization} 节）。
\end{itemize}

\subsection{实现目标}

围绕"清除率 $100\%$ 且耗时最短"，本方案的核心问题是：\textbf{在未知源位置、未知波束
方向、波束半径 $R\in[1000,1500]$ m 的前提下，用最少的测量位置与最省的访问路线，使
机器狗对\textbf{任意}定向源至少从正面命中一次}（检测保证），并在听到后完成交会定位与
清除（任务完成）。全文的检测布局在第~\ref{sec:detect} 节给出并验证，定位清除策略在第
\ref{sec:localization} 节给出，仿真验证与时间最优结论在第~\ref{sec:verify} 节给出。